\newpage

\section{Preventivo}
\subsection{Introduzione}
La presente sezione definisce il preventivo iniziale delle ore produttive e dei costi di realizzazione del progetto \textbf{Sonar}, in coerenza con quanto dichiarato dal gruppo \textbf{Git Happens} nella \textit{Dichiarazione impegni} presentata in fase di candidatura e nel rispetto delle tariffe orarie stabilite dal \textit{Regolamento del progetto didattico}. \\
\\
Ciascuno dei sette componenti del team prevede un impegno massimo di \textbf{93 ore produttive}, per un tetto complessivo di \textbf{651 ore produttive} e un preventivo economico stimato di \textbf{\texteuro{} 12.845,00}, assunto dal gruppo come limite massimo di spesa per l'intero progetto. \\
\\
I valori riportati nei paragrafi seguenti costituiscono la baseline di riferimento rispetto alla quale verranno confrontati i consuntivi di periodo al termine di ogni sprint, al fine di monitorare tempestivamente eventuali scostamenti e mantenere il progetto entro i vincoli di tempo e budget concordati.

\subsection{Costo orario per ruolo}
In conformità alle tariffe stabilite dal \textit{Regolamento del progetto didattico}, la seguente tabella riporta il costo orario associato a ciascuno dei sei ruoli di progetto.

\vspace{2mm}
\renewcommand{\arraystretch}{1.3}
\noindent
\begin{tabularx}{\textwidth}{|X c|}
\hline
\textbf{Ruolo} & \textbf{Costo orario (\texteuro{}/h)} \\
\hline
Responsabile & 30,00 \\
Amministratore & 20,00 \\
Analista & 25,00 \\
Progettista & 25,00 \\
Programmatore & 15,00 \\
Verificatore & 15,00 \\
\hline
\end{tabularx}

\newpage
\subsection{Prospetto orario per componente}
La seguente tabella riporta la ripartizione preventiva del monte orario massimo individuale (93 ore produttive) tra i sei ruoli di progetto per ciascuno dei sette componenti del gruppo \textbf{Git Happens}.

\vspace{2mm}
\renewcommand{\arraystretch}{1.3}
\noindent
{\setlength{\tabcolsep}{4.5pt}%
\begin{tabularx}{\textwidth}{|>{\raggedright\arraybackslash}X c c c c c c | c|}
\hline
\textbf{Componente} & \textbf{Resp.} & \textbf{Amm.} & \textbf{Anal.} & \textbf{Progett.} & \textbf{Programm.} & \textbf{Verif.} & \textbf{Totale} \\
\hline
Filippo Lissandrin & 7 & 7 & 10 & 20 & 24 & 25 & \textbf{93} \\
Cristiano Antoniazzo & 7 & 7 & 10 & 20 & 24 & 25 & \textbf{93} \\
Tommaso Caon & 7 & 7 & 10 & 20 & 24 & 25 & \textbf{93} \\
Gabriele Gambarro & 7 & 7 & 10 & 20 & 24 & 25 & \textbf{93} \\
Gianmaria Grinovero & 7 & 7 & 10 & 20 & 24 & 25 & \textbf{93} \\
Diana Grosselle & 7 & 7 & 10 & 20 & 24 & 25 & \textbf{93} \\
Antonio Hai Xu Tang & 7 & 7 & 10 & 20 & 24 & 25 & \textbf{93} \\
\hline
\textbf{TOTALE ORE} & \textbf{49} & \textbf{49} & \textbf{70} & \textbf{140} & \textbf{168} & \textbf{175} & \textbf{651} \\
\hline
\end{tabularx}}

\vspace{3mm}
\noindent
La suddivisione identica per tutti i componenti garantisce un carico di lavoro equo e una conoscenza trasversale dell'intero sistema, in linea con la rotazione dei ruoli adottata dal team. Il peso maggiore assegnato a \textbf{Verificatore} (175 ore), \textbf{Programmatore} (168 ore) e \textbf{Progettista} (140 ore) riflette i criteri architetturali e di qualità richiesti dal capitolato \textbf{C8 - Sonar}, mentre ai ruoli di \textbf{Analista} (70 ore), \textbf{Responsabile} e \textbf{Amministratore} (49 ore ciascuno) è allocata una quota coerente con le attività di analisi iniziale, coordinamento e gestione infrastrutturale.

\subsection{Totale ore rendicontate}
La tabella seguente definisce il modello di riferimento per il tracciamento delle ore rendicontate rispetto al monte ore preventivato per ciascun ruolo. Tale schema viene utilizzato per consolidare i dati consuntivi raccolti al termine di ogni sprint nella sezione \textit{Consuntivo di periodo} (Sezione~5), consentendo di verificare in modo immediato lo scostamento rispetto alla pianificazione iniziale e le ore residue disponibili.

\vspace{2mm}
\renewcommand{\arraystretch}{1.3}
\noindent
\begin{tabularx}{\textwidth}{|X c c c|}
\hline
\textbf{Ruolo} & \textbf{Ore preventivate} & \textbf{Ore rendicontate} & \textbf{Ore residue} \\
\hline
Responsabile & 49 & -- & -- \\
Amministratore & 49 & -- & -- \\
Analista & 70 & -- & -- \\
Progettista & 140 & -- & -- \\
Programmatore & 168 & -- & -- \\
Verificatore & 175 & -- & -- \\
\hline
\textbf{TOTALE} & \textbf{651} & \textbf{--} & \textbf{--} \\
\hline
\end{tabularx}

\newpage
\subsection{Totale complessivo}
Combinando le tariffe orarie definite nella Sezione~4.2 con il monte ore preventivato per ruolo nella Sezione~4.3, la tabella seguente sintetizza il costo stimato per ciascun ruolo, la relativa incidenza percentuale sul budget e il totale complessivo di progetto.

\vspace{2mm}
\renewcommand{\arraystretch}{1.3}
\noindent
\begin{tabularx}{\textwidth}{|X c c c c|}
\hline
\textbf{Ruolo} & \textbf{Costo orario (\texteuro{}/h)} & \textbf{Ore totali} & \textbf{Costo totale (\texteuro{})} & \textbf{Incidenza} \\
\hline
Responsabile & 30,00 & 49 & 1.470,00 & 11,44\% \\
Amministratore & 20,00 & 49 & 980,00 & 7,63\% \\
Analista & 25,00 & 70 & 1.750,00 & 13,62\% \\
Progettista & 25,00 & 140 & 3.500,00 & 27,25\% \\
Programmatore & 15,00 & 168 & 2.520,00 & 19,62\% \\
Verificatore & 15,00 & 175 & 2.625,00 & 20,44\% \\
\hline
\textbf{TOTALE} & \textbf{--} & \textbf{651} & \textbf{12.845,00} & \textbf{100,00\%} \\
\hline
\end{tabularx}

\vspace{3mm}
\noindent
Il preventivo economico complessivo per la realizzazione del progetto \textbf{Sonar} ammonta a \textbf{\texteuro{} 12.845,00} (a fronte di un tetto massimo di \textbf{651 ore produttive}), assunto dal gruppo \textbf{Git Happens} come limite di budget non superabile. Qualora il monitoraggio a consuntivo evidenzi il rischio di sforamento, il team interverrà rimodulando la distribuzione delle attività o negoziando con il proponente una riduzione concordata dell'ambito (a partire dai requisiti desiderabili o opzionali), senza incrementare il costo complessivo preventivato.

